\section{Vectors}

Vectors possess both

\begin{itemize}
	\item magnitude,
	\item direction.
\end{itemize}

Examples include

\begin{itemize}
	\item displacement,
	\item velocity,
	\item acceleration,
	\item force,
	\item momentum,
	\item electric field,
	\item magnetic field,
	\item probability current.
\end{itemize}

Vectors are usually written in boldface,

\[
\mathbf A,
\]

or with an arrow,

\[
\vec A.
\]

%%%%%%%%%%%%%%%%%%%%%%%%%%%%%%%%%%%%%%%%%%%%%%%%%%%%%%


\begin{bookfigure}{A vector resolved into Cartesian components. The same
geometric vector can be described by different components after a change of
coordinates.}
\begin{tikzpicture}[scale=1.05]
  \draw[physics axis] (-0.2,0) -- (5.6,0) node[right] {$x$};
  \draw[physics axis] (0,-0.2) -- (0,4.2) node[above] {$y$};
  \draw[physics vector] (0,0) -- (4.4,3.15)
    node[above right,visual label] {$\mathbf A$};
  \draw[construction line] (4.4,0) -- (4.4,3.15);
  \draw[construction line] (0,3.15) -- (4.4,3.15);
  \draw[physics vector secondary] (0,-0.25) -- (4.4,-0.25)
    node[midway,below] {$A_x$};
  \draw[physics vector secondary] (-0.25,0) -- (-0.25,3.15)
    node[midway,left] {$A_y$};
  \draw[visualgold,line width=0.9pt,-{Latex[length=1.5mm]}]
    (0.9,0) arc[start angle=0,end angle=35.6,radius=0.9];
  \node[text=visualgold,font=\small] at (1.12,0.36) {$\theta$};
\end{tikzpicture}
\end{bookfigure}

\begin{visualguidebox}
The diagonal arrow is the vector itself. The horizontal and vertical arrows are
its components relative to the selected Cartesian basis.
\end{visualguidebox}
